\subsubsection{Export Bulk Applicant Lists}
The use case referred to as Export Bulk Applicant Lists gives the registrar a tool for exporting full data of applicants for management purposes. After signing into the application using the required credentials, the user gets access to the main applicant list. The system gives the ability to export applicant information in files. It will assist in record keeping and reporting. It is mostly used when closing off the recruitment process for the academic year.
\begin{table}[H]
    \centering
    \small
    \renewcommand{\arraystretch}{1.4}
    \caption{Use Case Table: Export Bulk Applicant Lists (UC-5.13)} % ক্যাপশন উপরে আনা হয়েছে
    \label{tab:export-bulk-applicant-registrar}
    \begin{tabularx}{\textwidth}{|l|X|}
        \hline
        \rowcolor[gray]{0.9} {Characteristics} & {Description} \\ \hline
        Use Case ID & UC-5.13 \\ \hline
        Use Case Name & Export Bulk Applicant Lists \\ \hline
        Primary Actors & Registrar \\ \hline
        Goal & Export comprehensive lists of all applicants for high-level administrative reporting. \\ \hline
        Precondition & User must be logged into the system with Registrar-level administrative privileges. \\ \hline
        Scenario & 1. Login to the system. 2. Go to the master list of applicants. 3. Click the 'Export' button to create files for all applicants. \\ \hline
        Exception & User does not have required role; timeout on server during processing of bulk action. \\ \hline
        Frequency of Use & High frequency towards the end of admission cycle. \\ \hline
    \end{tabularx}
\end{table}
\FloatBarrier